% Generated from 2/2-3.md; technical content preserved.
\exercise{۲.۳}
در \textbf{\lr{Exercise 2.3}} صورت مسئله بیان می‌کند که اگر در مدار \textbf{\lr{Figure 2.32}} اثر \textbf{\lr{Body Effect}} لحاظ شود، آیا جریان دو ترانزیستور سری (\(I_{DS2}\)) همچنان برابر، بیشتر یا کمتر از جریان ترانزیستور منفرد (\(I_{DS1}\)) خواهد بود؟

\subsection*{۱. پاسخ مستقیم}
با در نظر گرفتن \textbf{\lr{Body Effect}}، جریان دو ترانزیستور سری کمتر از ترانزیستور منفرد خواهد شد:
\[\mathbf{I_{DS2} < I_{DS1}} \quad \text{\rl{یا}} \quad \mathbf{I_{DS1} > I_{DS2}}\]

\subsection*{۲. تحلیل فیزیکی و مداری}

\subsubsection*{الف) ترانزیستور منفرد در مدار (\lr{a})}
\begin{itemize}
\item در ترانزیستور منفرد با طول \(2L\)، سورس (\lr{Source}) به زمین (\lr{GND}) وصل است.
\item بدنه (\lr{Substrate/Body}) ترانزیستور \lr{nMOS} نیز به پتانسیل زمین وصل می‌شود.
\item بنابراین ولتاژ سورس-بدنه برابر صفر است:
\[V_{SB} = V_S - V_B = 0 - 0 = 0\]
\item در نتیجه ولتاژ آستانه (\lr{Threshold Voltage}) تغییری نکرده و دقیقاً برابر با ولتاژ آستانه نامی صفر بایاس (\(V_{t0}\)) باقی می‌ماند:
\[V_t = V_{t0}\]

\end{itemize}
\subsubsection*{ب) ترانزیستورهای سری در مدار (\lr{b})}
در ساختار سری دو ترانزیستور با ابعاد \(W/L\):
\begin{enumerate}
\item \textbf{ترانزیستور پایینی (\lr{Bottom Transistor}):}
\begin{itemize}
\item سورس این ترانزیستور به زمین متصل است؛ پس \(V_{SB,bottom} = 0\) بوده و ولتاژ آستانه آن بدون تغییر برابر با \(V_{t0}\) است.

\end{itemize}
\item \textbf{ترانزیستور بالایی (\lr{Top Transistor}):}
\begin{itemize}
\item سورس ترانزیستور بالایی به گره میانی با پتانسیل \(V_1\) متصل است.
\item هنگام برقراری جریان، پتانسیل گره میانی از صفر بیشتر است (\(V_1 > 0\)).
\item از آن‌جا که بدنه هر دو ترانزیستور \lr{nMOS} بر روی یک سابستریت مشترک \lr{p-type} متصل به زمین (\lr{GND}) قرار دارد، ولتاژ سورس-بدنه ترانزیستور بالایی مثبت می‌شود:
\[V_{SB,top} = V_1 - 0 = V_1 > 0\]

\end{itemize}
\end{enumerate}

\subsection*{۳. تغییرات ولتاژ آستانه و کاهش جریان}
رابطه تحلیلی \textbf{\lr{Body Effect}} برای ولتاژ آستانه به صورت زیر است:
\[V_t = V_{t0} + \gamma \left( \sqrt{\phi_s + V_{SB}} - \sqrt{\phi_s} \right)\]
یا در تقریب خطی برای مقادیر کوچک ولتاژ سورس:
\[V_t \approx V_{t0} + k_\gamma V_{SB}\]

از آن‌جا که برای ترانزیستور بالایی \(V_{SB,top} = V_1 > 0\) است:
\begin{itemize}
\item بار ناحیه تخلیه زیر گیت (\lr{Depletion Charge}) افزایش یافته و جهت ایجاد لایه وارونگی (\lr{Inversion Layer}) به ولتاژ گیت بیشتری نیاز است.
\item در نتیجه ولتاژ آستانه ترانزیستور بالایی افزایش می‌یابد:
\[V_{t,top} > V_{t0}\]
\item با افزایش \(V_{t,top}\)، ولتاژ مؤثر گیت (\lr{Gate Overdrive Voltage}) یعنی \(V_{GT,top} = V_{GS,top} - V_{t,top}\) کاهش می‌یابد.
\item کاهش \(V_{GT}\) باعث کاهش مستقیم جریان قابل‌انتقال توسط ترانزیستور بالایی و افزایش مقاومت حالت روشن آن می‌شود.
\item چون این دو ترانزیستور سری هستند، کل جریان شاخه توسط ترانزیستور ضعیف‌تر محدود شده و نقطه کار گره \(V_1\) به گونه‌ای تنظیم می‌شود که جریان شاخه در نهایت کمتر از حالت ایده‌آل بدون اثر بدنه شود.

\end{itemize}
بنابراین:
\[I_{DS2} < I_{DS1}\]

\subsection*{۴. پیامد این پدیده در طراحی \lr{VLSI}}
این تمرین مستقیماً دلیل افت کارایی در \textbf{\lr{Series Stacks}} (مانند شبکه \lr{Pull-Down} در گیت‌های منطقی \lr{NAND} با تعداد ورودی زیاد) را مدل می‌کند؛ هر چه ترانزیستور در استک بالاتر قرار گیرد، \(V_{SB}\) آن بزرگ‌تر شده و با افزایش \(V_t\)، جریان هدایت شاخه بیشتر تضعیف می‌شود.
